\documentclass[11pt]{article}

\usepackage[margin=1in]{geometry}
\usepackage{amsmath,amssymb,amsthm}
\usepackage{enumitem}
\usepackage{graphicx}
\usepackage[colorlinks=true,linkcolor=blue,urlcolor=blue,citecolor=blue]{hyperref}
\emergencystretch=2em

\newtheorem{theorem}{Theorem}
\newtheorem{corollary}[theorem]{Corollary}
\theoremstyle{remark}
\newtheorem{remark}[theorem]{Remark}

\newcommand{\Dc}{\mathcal D_c^+}
\newcommand{\Delp}{\Delta^+}
\newcommand{\Lmax}{M^*}
\newcommand{\tens}{\mathbin{\otimes}}

\title{Continuity of the t-Norm Is Necessary\\
in Two Broad t-Conorm Regimes}
\author{Math discovery candidate partial result}
\date{August 17, 2026}

\begin{document}
\maketitle

\section{Status and source question}

Lai and Luo prove that if a t-norm \(T\) is continuous and a
right-continuous t-conorm \(L\) satisfies their strictness condition
\((LS)\), then the tensor product \(L\tens T\) preserves the continuous
distance distribution functions \(\Dc\).  Remark 4.7 asks whether continuity
of \(T\) is necessary \cite{source}.  The question appears on printed page 15:

\begin{center}
\includegraphics[width=\textwidth]{figures/open_problem_crop.png}
\end{center}

This packet proves necessity exactly for the maximum t-conorm.  It also
excludes every left-continuous t-norm having a nilpotent-minimum ordinal-sum
component when \(L\) is addition, or more generally has an additive
generator.  The arbitrary-\(L\) question remains open.

\section{Definitions and proof intuition}

A distance distribution function (d.d.f.) is an increasing, left-continuous
map \(f:[0,\infty]\to[0,1]\) with \(f(0)=0\) and \(f(\infty)=1\).
Write \(\Dc\) for the continuous d.d.f.s.  For finite \(x\),
\[
 (L\tens T)(f,g)(x)
   =\sup_{L(r,s)<x}T(f(r),g(s)).                              \tag{1}
\]

For \(L=\Lmax\), where \(\Lmax(r,s)=\max\{r,s\}\), the source observes
that (1) collapses to the pointwise composition
\(T(f(x),g(x))\).  Continuous d.d.f.s can pass through any prescribed point
of the unit square, so preservation of continuity tests all northeast limits
of \(T\).  Left continuity already controls southwest limits; monotonicity
then squeezes every other approach.

For addition, a nilpotent-minimum component has a threshold: below it the
t-norm stays at the component floor, but immediately above it the value
jumps to the midpoint.  A piecewise-linear d.d.f. turns that algebraic
threshold into a literal spatial discontinuity of (1).

\section{The exact maximum-t-conorm subcase}

\begin{theorem}[Maximum t-conorm]
\label{thm:max}
Let \(T\) be a left-continuous t-norm.  Then
\[
       (\Dc,\Lmax\tens T)\text{ is closed under }\Lmax\tens T
       \quad\Longleftrightarrow\quad
       T\text{ is continuous on }[0,1]^2.
\]
\end{theorem}

\begin{proof}
The source paper computes, for every finite \(x\),
\[
 (\Lmax\tens T)(f,g)(x)
 =\sup_{r,s<x}T(f(r),g(s))=T(f(x),g(x)).                     \tag{2}
\]
If \(T\) is continuous, (2) is continuous on the finite half-line for every
\(f,g\in\Dc\).  At infinity, \(f(x),g(x)\to1\), so the limit is
\(T(1,1)=1\); hence the output lies in \(\Dc\).

Conversely, assume the stated closure.  Fix \(p,q\in(0,1)\), and define
continuous d.d.f.s
\[
 f_p(x)=
 \begin{cases}
   px,&0\le x\le1,\\
   p+(1-p)(x-1),&1\le x\le2,\\
   1,&x\ge2,
 \end{cases}
 \qquad
 g_q(x)=
 \begin{cases}
   qx,&0\le x\le1,\\
   q+(1-q)(x-1),&1\le x\le2,\\
   1,&x\ge2.
 \end{cases}                                                   \tag{3}
\]
By closure and (2), \(x\mapsto T(f_p(x),g_q(x))\) is continuous at \(1\).
Therefore
\[
 \lim_{t\downarrow0}T\bigl(p+(1-p)t,q+(1-q)t\bigr)=T(p,q).    \tag{4}
\]
Because \(T\) is increasing in each coordinate, any path approaching
\((p,q)\) strictly from the northeast is squeezed between two
reparameterizations of the path in (4).  Thus (4) is the full northeast
limit of \(T\) at \((p,q)\).  The assumed left continuity gives
\[
 \lim_{t\downarrow0}T(p-t,q-t)=T(p,q),                         \tag{5}
\]
with truncated coordinates near the boundary.

Now let \((u_n,v_n)\to(p,q)\).  For every sufficiently small
\(\delta>0\), eventually
\[
 T(p-\delta,q-\delta)
 \le T(u_n,v_n)
 \le T(p+\delta,q+\delta),                                    \tag{6}
\]
again truncating at \(0\) and \(1\).  Equations (4)--(5), followed by
\(\delta\downarrow0\), squeeze \(T(u_n,v_n)\) to \(T(p,q)\).
If \(p=0\) or \(q=0\), continuity follows directly from
\(0\le T(u,v)\le\min\{u,v\}\).  If \(p=1\) or \(q=1\), use the
left-continuity squeeze in that coordinate and \(T(1,v)=v\) in the other.
Thus \(T\) is continuous everywhere.
\end{proof}

\section{An explicit ordinal-sum obstruction}

Let the nilpotent-minimum t-norm be
\[
 N(u,v)=
 \begin{cases}
  0,&u+v\le1,\\
  \min\{u,v\},&u+v>1.
 \end{cases}                                                    \tag{7}
\]
We say that a t-norm \(T\) has a nilpotent-minimum ordinal-sum component
on \([a,b]\), \(0\le a<b\le1\), if on that square it is
\[
 T(u,v)=a+(b-a)N\left(\frac{u-a}{b-a},\frac{v-a}{b-a}\right),   \tag{8}
\]
and the usual ordinal-sum rule applies outside a common component (so the
value is the minimum).

\begin{theorem}[Addition]
\label{thm:add}
If a left-continuous t-norm \(T\) has a nilpotent-minimum ordinal-sum
component, then \(\Dc\) is not closed under \(+\tens T\).
\end{theorem}

\begin{proof}
For the component \([a,b]\), define the continuous d.d.f.
\[
 F(x)=
 \begin{cases}
   ax,&0\le x\le1,\\
   a+(b-a)(x-1),&1\le x\le2,\\
   b+(1-b)(x-2),&2\le x\le3,\\
   1,&x\ge3.
 \end{cases}                                                    \tag{9}
\]
Set \(H=(+\tens T)(F,F)\).  We claim \(H(3)=a\).  Indeed, if
\(r+s<3\) and both \(r,s\in[1,2]\), their normalized coordinates in (8)
sum to less than \(1\), so \(T(F(r),F(s))=a\).  If one spatial coordinate
is below \(1\), one t-norm input is at most \(a\).  If one is above \(2\),
the other must be below \(1\), and the ordinal-sum rule again gives a value
at most \(a\).  Hence the supremum in (1) is at most \(a\), while pairs in
\([1,2]^2\) with sum below \(3\) show that it is exactly \(a\).

For every \(x>3\), choose \(\varepsilon>0\) with
\(2(3/2+\varepsilon)<x\).  Both spatial coordinates then lie in the
component, their normalized sum is greater than \(1\), and
\[
 H(x)\ge
 T\bigl(F(3/2+\varepsilon),F(3/2+\varepsilon)\bigr)
 =a+(b-a)(1/2+\varepsilon).                                   \tag{10}
\]
Consequently
\(\liminf_{x\downarrow3}H(x)\ge a+(b-a)/2>a=H(3)\).
Thus \(H\notin\Dc\).
\end{proof}

For \(a=0,b=1\), (9) can be shortened to \(F(x)=\min\{x,1\}\); then
\((+\tens N)(F,F)(1)=0\), while its right limit is at least \(1/2\).

\begin{corollary}[Additive-generator t-conorms]
\label{cor:generator}
Suppose a t-conorm on the extended half-line has the form
\[
 L(r,s)=\phi^{-1}\bigl(\phi(r)+\phi(s)\bigr),                  \tag{11}
\]
where \(\phi:[0,\infty]\to[0,\infty]\) is an increasing homeomorphism.
If \(T\) has a nilpotent-minimum ordinal-sum component, then \(\Dc\) is not
closed under \(L\tens T\).
\end{corollary}

\begin{proof}
Let \(F\) be (9) and put \(f=F\circ\phi\).  This is a continuous d.d.f.
For \(X<\infty\), the inequality
\(L(r,s)<\phi^{-1}(X)\) is equivalent to
\(\phi(r)+\phi(s)<X\).  Therefore
\[
 (L\tens T)(f,f)(\phi^{-1}(X))=(+\tens T)(F,F)(X).
\]
The jump at \(X=3\) from Theorem~\ref{thm:add} is transported to
\(\phi^{-1}(3)\).
\end{proof}

\section{What remains open}

Theorem~\ref{thm:max} completely answers Remark 4.7 for one noncancellative
t-conorm satisfying \((LS)\).  Theorem~\ref{thm:add} and
Corollary~\ref{cor:generator} rule out a canonical, broad discontinuous
family for all cancellative t-conorms with additive generators.  They do not
cover every discontinuous left-continuous t-norm or every right-continuous
t-conorm satisfying \((LS)\).

The source's quasi-inverse formula suggests the missing general step.  A
discontinuity changes a strict superlevel set of \(T\), but one would need to
prove that this change is detected by
\[
 \inf_{T(p,q)>y}L(u(p),v(q))                                   \tag{12}
\]
for some strictly increasing, infimum-preserving \(u,v\).  For a general
\(L\), other points of the superlevel set may keep the infimum in (12)
unchanged.  Eight focused upgrade attempts did not remove this
order-separation obstruction.

\section{Verification and novelty status}

The maximum proof was audited at interior and boundary points.  The ordinal
witness was checked over all spatial cases \(r+s<3\), and the generator
corollary is an exact change of variables.  The accompanying verification
report records those checks.

Bounded searches through August 17, 2026 by arXiv id, exact title, exact
remark wording, and the relevant t-norm/t-conorm phrases found the final
journal publication but no later answer to Remark 4.7.  The source already
contains identity (2); the exact continuity equivalence and the explicit
ordinal-sum obstruction were not located.  This is therefore a candidate
new partial result, not a certified novelty claim.

\begin{thebibliography}{9}
\bibitem{source}
H. Lai and Q. Luo,
\emph{Triangle functions generated by tensor products of quantales},
arXiv:2411.11876 (2024); Fuzzy Sets and Systems 508 (2026), 109655.

\bibitem{tnorms}
E. P. Klement, R. Mesiar, and E. Pap,
\emph{Triangular Norms}, Kluwer Academic Publishers, 2000.

\bibitem{pm}
B. Schweizer and A. Sklar,
\emph{Probabilistic Metric Spaces}, Dover, 2005.
\end{thebibliography}

\end{document}
